% ============================================================
% Reading notes on
% "A Polynomial-Time Algorithm for 1/2 Well-Supported Nash
%  Equilibria in Bimatrix Games" (Fearnley & Savani)
% Translated and adapted from the Chinese notes (WSNE.md / WSNE.tex).
% ============================================================

\section{Preliminaries}

For a bimatrix game, after normalizing the payoff matrices, we can define $\varepsilon$-NE and $\varepsilon$-WSNE:

\begin{itemize}
  \item \textbf{$\varepsilon$-NE ($\varepsilon$-Nash Equilibrium)} \cite{nashtheorem51}: Formally, a mixed strategy profile $(\mathbf{x}, \mathbf{y})$ is called an $\varepsilon$-NE if and only if:
  \begin{align*}
    \mathcal{R}(\mathbf{x}, \mathbf{y}) &\ge \max_{i \in [n]} \mathcal{R}(e_i, \mathbf{y}) - \varepsilon \\
    \mathcal{C}(\mathbf{x}, \mathbf{y}) &\ge \max_{j \in [n]} \mathcal{C}(\mathbf{x}, e_j) - \varepsilon
  \end{align*}
  That is, the maximum gain of any player from a unilateral deviation is at most $\varepsilon$.

  \item \textbf{$\varepsilon$-WSNE ($\varepsilon$-Well-Supported Nash Equilibrium)}: Formally, a mixed strategy profile $(\mathbf{x}, \mathbf{y})$ is called an $\varepsilon$-WSNE if and only if:
  \begin{align*}
    \mathcal{R}(e_i, \mathbf{y}) &\ge \max_{k \in [n]} \mathcal{R}(e_k, \mathbf{y}) - \varepsilon && \text{for all } i \in \operatorname{supp}(\mathbf{x}) \\
    \mathcal{C}(\mathbf{x}, e_j) &\ge \max_{k \in [n]} \mathcal{C}(\mathbf{x}, e_k) - \varepsilon && \text{for all } j \in \operatorname{supp}(\mathbf{y})
  \end{align*}
  That is, every pure strategy in the support of the current strategy is an $\varepsilon$-best response, i.e., the gain from deviating from it is at most $\varepsilon$.
\end{itemize}

It is easy to see that every $\varepsilon$-WSNE is an $\varepsilon$-NE (multiplying the inequalities for all pure strategies in the support by their probabilities and summing them up yields the two inequalities of an $\varepsilon$-NE). Intuitively, $\varepsilon$-WSNE avoids the situation where a player assigns small probabilities to some ``bad strategies'' to gain payoff, since these bad strategies may become the bottleneck in practice.

The algorithms below will repeatedly use the LMM lemma and its corollaries for finding approximate equilibria in general bimatrix games:

\begin{lemmaBox}{Lipton--Markakis--Mehta Small-Support Lemma [LMM03]}
  Let $(\mathbf{x}, \mathbf{y})$ be a NE of an $n \times n$ bimatrix game $(R, C)$. For any $\varepsilon > 0$ and $k \ge \dfrac{12 \ln n}{\varepsilon^2}$, there exists a $k$-uniform strategy profile $(\mathbf{x}^*, \mathbf{y}^*)$ such that $(\mathbf{x}^*, \mathbf{y}^*)$ is an $\varepsilon$-NE and
  \begin{align*}
    |\mathcal{R}(\mathbf{x}^*, \mathbf{y}^*) - \mathcal{R}(\mathbf{x}, \mathbf{y})| &\le \varepsilon \\
    |\mathcal{C}(\mathbf{x}^*, \mathbf{y}^*) - \mathcal{C}(\mathbf{x}, \mathbf{y})| &\le \varepsilon
  \end{align*}

  In \cite{lmm03}, a probabilistic method was used to prove that there exists a sampling satisfying the above properties:

  Sample $k$ pure strategies from $\mathbf{x}$ independently and with replacement to form a multiset $A$ (note that the sampling is according to the distribution $\mathbf{x}$), and similarly sample $k$ pure strategies from $\mathbf{y}$ to form a multiset $B$. Assign probability $1/k$ to each pure strategy in $A$ and $B$ to obtain the mixed strategies $\mathbf{x}^*, \mathbf{y}^*$. Clearly, if a pure strategy appears $t$ times in $A$, its probability is $t / k$.

  Define the following events:
  \begin{align*}
    \phi_1 &= \{ |\mathcal{R}(\mathbf{x}^*, \mathbf{y}^*) - \mathcal{R}(\mathbf{x}, \mathbf{y})| < \varepsilon/2 \} \\
    \phi_2 &= \{ |\mathcal{C}(\mathbf{x}^*, \mathbf{y}^*) - \mathcal{C}(\mathbf{x}, \mathbf{y})| < \varepsilon/2 \} \\
    \pi_{1, i} &= \{ \mathcal{R}(e_i, \mathbf{y}^*) < \mathcal{R}(\mathbf{x}^*, \mathbf{y}^*) + \varepsilon \} \quad i \in [n] \\
    \pi_{2, j} &= \{ \mathcal{C}(\mathbf{x}^*, e_j) < \mathcal{C}(\mathbf{x}^*, \mathbf{y}^*) + \varepsilon \} \quad j \in [n] \\
    \text{GOOD} &= \phi_1 \cap \phi_2 \bigcap_{i \in [n]} \pi_{1, i} \bigcap_{j \in [n]} \pi_{2, j}
  \end{align*}
  The goal is to prove $\mathbb{P}(\text{GOOD}) > 0$. Note that the deviation from the NE is shrunk to $\varepsilon / 2$ here, because we will later need it to bound $\overline{\pi_{1, i}}$ and $\overline{\pi_{2, i}}$.

  Consider the $\phi$ events. To bound $\overline{\phi_1}$ and $\overline{\phi_2}$, define the following events:
  \begin{align*}
    \phi_{1, a} &= \{ |\mathcal{R}(\mathbf{x}^*, \mathbf{y}) - \mathcal{R}(\mathbf{x}, \mathbf{y})| < \varepsilon/4 \} \\
    \phi_{1, b} &= \{ |\mathcal{R}(\mathbf{x}^*, \mathbf{y}^*) - \mathcal{R}(\mathbf{x}^*, \mathbf{y})| < \varepsilon/4 \} \\
    \phi_{2, a} &= \{ |\mathcal{C}(\mathbf{x}, \mathbf{y}^*) - \mathcal{C}(\mathbf{x}, \mathbf{y})| < \varepsilon/4 \} \\
    \phi_{2, b} &= \{ |\mathcal{C}(\mathbf{x}^*, \mathbf{y}^*) - \mathcal{C}(\mathbf{x}, \mathbf{y}^*)| < \varepsilon/4 \}
  \end{align*}
  Then $\phi_{1, a} \cap \phi_{1, b} \subseteq \phi_1$, and similarly for $\phi_2$.

  Since every pure strategy in $\operatorname{supp}(\mathbf{x}^*) \subseteq \operatorname{supp}(\mathbf{x})$ is a best response to $\mathbf{y}$, its expected payoff equals $\mathcal{R}(\mathbf{x}, \mathbf{y})$. Viewing the payoffs of the $k$ pure strategies in the multiset $A$ as random variables, each takes values in $(0, 1]$ with expectation $\mathcal{R}(\mathbf{x}, \mathbf{y})$. By Hoeffding's inequality:

  \begin{lemmaBox}{Hoeffding's Inequality}
    For $n$ independent random variables $X_i \in [l_i, r_i]$ with empirical mean $\overline{X} = \dfrac{\sum_{i} X_i}{n}$,
    \[
      \forall \varepsilon > 0,\ \mathbb{P}(\overline{X} - \mathbb{E}(\overline{X}) \ge \varepsilon) \le \exp\left( -\dfrac{2n^2\varepsilon^2}{\sum_{i} (r_i - l_i)^2} \right)
    \]
  \end{lemmaBox}

  we get:
  \begin{align*}
    \mathbb{P}(\mathcal{R}(\mathbf{x}^*, \mathbf{y}) - \mathcal{R}(\mathbf{x}, \mathbf{y}) \ge \varepsilon / 4) &\le \exp(-k\varepsilon^2 / 8) \\
    \mathbb{P}(\mathcal{R}(\mathbf{x}, \mathbf{y}) - \mathcal{R}(\mathbf{x}^*, \mathbf{y}) \ge \varepsilon / 4) &\le \exp(-k\varepsilon^2 / 8)
  \end{align*}
  i.e., $\mathbb{P}(\overline{\phi_{1, a}}) \le 2\exp(-k\varepsilon^2 / 8)$; similarly $\mathbb{P}(\overline{\phi_{1, b}}) \le 2\exp(-k\varepsilon^2 / 8)$, and the same holds for $\phi_2$. Therefore,
  \begin{align*}
    \mathbb{P}(\overline{\phi_1}) &\le 4\exp(-k\varepsilon^2 / 8) \\
    \mathbb{P}(\overline{\phi_2}) &\le 4\exp(-k\varepsilon^2 / 8)
  \end{align*}

  Consider the $\pi$ events. Define the following events:
  \begin{align*}
    \psi_{1, i} &= \{ \mathcal{R}(e_i, \mathbf{y}^*) < \mathcal{R}(e_i, \mathbf{y}) + \varepsilon / 2 \} \quad i \in [n] \\
    \psi_{2, j} &= \{ \mathcal{C}(\mathbf{x}^*, e_j) < \mathcal{C}(\mathbf{x}, e_j) + \varepsilon / 2 \} \quad j \in [n]
  \end{align*}
  Earlier we set $\phi_1, \phi_2$ to tighter bounds, which can now be combined with the new events. For example, on $\phi_1 \cap \psi_{1, i}$:
  \begin{align*}
    \mathcal{R}(e_i, \mathbf{y}^*) &< \mathcal{R}(e_i, \mathbf{y}) + \varepsilon / 2 \le \mathcal{R}(\mathbf{x}, \mathbf{y}) + \varepsilon / 2 \tag{$(\mathbf{x}, \mathbf{y})$ is a NE} \\
    &\le \mathcal{R}(\mathbf{x}^*, \mathbf{y}^*) + \varepsilon
  \end{align*}
  This gives $\pi_{1, i}$, and $\pi_{2, j}$ follows similarly. Likewise, for the $k$ pure strategies in the multiset $B$, viewing their payoffs against $e_i$ as random variables with expectation $\mathcal{R}(e_i, \mathbf{y})$ and range $(0, 1]$, Hoeffding's inequality gives:
  \[
    \mathbb{P}(\mathcal{R}(e_i, \mathbf{y}^*) - \mathcal{R}(e_i, \mathbf{y}) \ge \varepsilon / 2) \le \exp(-k\varepsilon^2 / 2)
  \]
  and similarly for $\psi_{2, j}$. Combining all the above parts, we have
  \begin{align*}
    \mathbb{P}(\overline{\text{GOOD}}) &\le \mathbb{P}(\overline{\phi_1}) + \mathbb{P}(\overline{\phi_2}) + \sum_{i \in [n]} \mathbb{P}(\overline{\pi_{1, i}}) + \sum_{j \in [n]} \mathbb{P}(\overline{\pi_{2, i}}) \\
    &\le (8n + 8) \exp(-k\varepsilon^2 / 8) + 2n \exp(-k\varepsilon^2 / 2)
  \end{align*}
  Plugging in the lower bound on $k$ gives $\mathbb{P}(\overline{\text{GOOD}}) < 1$, hence $\mathbb{P}(\text{GOOD}) > 0$, which means there exists at least one sampling satisfying the above properties.
\end{lemmaBox}

\section{Main Results}

\begin{resultBox}{Main Conclusion}
  For any constant $\delta > 0$, there is a polynomial-time algorithm that finds a $(1/2 + \delta)$-WSNE in any $n \times n$ normalized bimatrix game $(R, C)$.
\end{resultBox}

Note that the exponent of the ``polynomial time'' depends on $\delta$; when $\delta$ is a constant, the running time is polynomial in $n$.

For games in which the payoff matrix is not given but only a payoff-query oracle is available, the paper also gives a query version:

\begin{resultBox}{Query Version}
  For any constants $\varepsilon, \delta > 0$, with $O\left( \frac{n \log n}{\varepsilon^4} \right)$ payoff queries, one can find, in polynomial time and with success probability $\left( 1 - n^{-\frac{1}{8}} \right)\left( 1 - \frac{2}{n} \right)^2$, a $(1/2 + 3\varepsilon + \delta)$-WSNE.
\end{resultBox}

In general, to compute a $(1/2 + \delta')$-WSNE, one can take $\varepsilon = \delta = \frac{\delta'}{4}$. Similarly, the polynomial time here also depends on $\varepsilon$ and $\delta$.

\section{Algorithm 1: Polynomial-Time \texorpdfstring{$(1/2 + \delta)$}{1/2 + δ}-WSNE Computation}

\subsection{Preprocessing}

\begin{enumerate}
  \item Compute a Nash equilibrium $(\mathbf{x}^*, \mathbf{y}^*)$ of the zero-sum game $(R, -R)$; this can be solved in polynomial time via linear programming.

  By the properties of zero-sum games, $\mathbf{x}^*$ is the row player's \textbf{maximin strategy}:
  \[
    \mathcal{R}(\mathbf{x}^*, \mathbf{y}^*) = \max_{\mathbf{x} \in \Delta^n} \min_{\mathbf{y} \in \Delta^n} \mathcal{R}(\mathbf{x}, \mathbf{y})
  \]
  Hence for every column strategy $\mathbf{y}$, $\mathcal{R}(\mathbf{x}^*, \mathbf{y}^*) \le \mathcal{R}(\mathbf{x}^*, \mathbf{y})$. \textbf{This property will be used several more times later.}

  \item Compute a Nash equilibrium $(\hat{\mathbf{x}}, \hat{\mathbf{y}})$ of the zero-sum game $(-C, C)$. Similarly, for every row strategy $\mathbf{x}$, $\mathcal{C}(\hat{\mathbf{x}}, \hat{\mathbf{y}}) \le \mathcal{C}(\mathbf{x}, \hat{\mathbf{y}})$.
\end{enumerate}

We then case-analyze the relationship between $\mathcal{R}(\mathbf{x}^*, \mathbf{y}^*)$ and $\mathcal{C}(\hat{\mathbf{x}}, \hat{\mathbf{y}})$. By symmetry, we only analyze the case \textbf{$\mathcal{R}(\mathbf{x}^*, \mathbf{y}^*) \ge \mathcal{C}(\hat{\mathbf{x}}, \hat{\mathbf{y}})$}:

\subsection{The \texorpdfstring{$\mathcal{R}(\mathbf{x}^*, \mathbf{y}^*) \ge \mathcal{C}(\hat{\mathbf{x}}, \hat{\mathbf{y}})$}{R(x*,y*) >= C(xh,yh)} branch}

\paragraph{Case 1(a): $\mathcal{R}(\mathbf{x}^*, \mathbf{y}^*) \le 1/2$}
In this case, we simply output $(\hat{\mathbf{x}}, \mathbf{y}^*)$. We show below that it is a $1/2$-WSNE:

\begin{itemize}
  \item Row player: since $(\mathbf{x}^*, \mathbf{y}^*)$ is a NE of $(R, -R)$, for every pure strategy $i$, $\mathcal{R}(e_i, \mathbf{y}^*) \le \mathcal{R}(\mathbf{x}^*, \mathbf{y}^*) \le 1/2$. Hence the row player's best-response payoff is $\le 1/2$, so every row strategy is a $1/2$-best response to $\mathbf{y}^*$;
  \item Column player: since $(\hat{\mathbf{x}}, \hat{\mathbf{y}})$ is a NE of $(-C, C)$ and the case condition gives $\mathcal{C}(\hat{\mathbf{x}}, \hat{\mathbf{y}}) \le \mathcal{R}(\mathbf{x}^*, \mathbf{y}^*) \le 1/2$, for every pure strategy $j$, $\mathcal{C}(\hat{\mathbf{x}}, e_j) \le \mathcal{C}(\hat{\mathbf{x}}, \hat{\mathbf{y}}) \le \mathcal{R}(\mathbf{x}^*, \mathbf{y}^*) \le 1/2$, so every column strategy is a $1/2$-best response to $\hat{\mathbf{x}}$.
\end{itemize}

Therefore, the strategy profile $(\hat{\mathbf{x}}, \mathbf{y}^*)$ is a $1/2$-WSNE.

\paragraph{Case 1(b): ``one high, one low''}
In this case, $\mathcal{R}(\mathbf{x}^*, \mathbf{y}^*) > 1/2$, and there \textbf{exists} a strategy $\mathbf{x}'$ such that:
\begin{itemize}
  \item $\operatorname{supp}(\mathbf{x}') \subseteq \operatorname{supp}(\mathbf{x}^*)$;
  \item $\mathcal{C}(\mathbf{x}', e_j) \le 1/2 \quad \text{for all } j \in [n]$.
\end{itemize}

We then output $(\mathbf{x}', \mathbf{y}^*)$.

\textbf{How to check existence?}
Solve the linear feasibility system:
\[
  \begin{cases}
    \sum_{i \in \operatorname{supp}(\mathbf{x}^*)} \mathbf{x}'(i) \cdot C_{ij} \le \frac12, & \forall j \in [n] \\[4pt]
    \mathbf{x}'(i) \ge 0, & \forall i \in \operatorname{supp}(\mathbf{x}^*) \\[4pt]
    \sum_{i \in \operatorname{supp}(\mathbf{x}^*)} \mathbf{x}'(i) = 1
  \end{cases}
\]
This linear program can be solved, or declared infeasible, in polynomial time.

We now show that $(\mathbf{x}', \mathbf{y}^*)$ is a $1/2$-WSNE:

\begin{itemize}
  \item Row player: since $\operatorname{supp}(\mathbf{x}') \subseteq \operatorname{supp}(\mathbf{x}^*)$ and every strategy in $\operatorname{supp}(\mathbf{x}^*)$ is a best response to $\mathbf{y}^*$, when the row player uses $\mathbf{x}'$, every pure strategy in the support is an exact best response, with payoff equal to $\mathcal{R}(\mathbf{x}^*, \mathbf{y}^*) > 1/2$. That is, the row player's strategy is a $0$-best response.
  \item Column player: by construction, $\mathcal{C}(\mathbf{x}', e_j) \le 1/2$ for all $j$, so the column player's best-response payoff is $\le 1/2$, and hence every column strategy is a $1/2$-best response to $\mathbf{x}'$.
\end{itemize}

\paragraph{Case 1(c): ``both high''}
In this case, neither of the previous two conditions holds; that is, we cannot bound a single player's best response by $\le 1/2$.

Define the \textbf{restricted game} $(R_{\mathbf{x}^*}, C_{\mathbf{x}^*})$: the row player may only use actions in $\operatorname{supp}(\mathbf{x}^*)$, while the column player may use all $n$ actions.

One can show that every Nash equilibrium $(\mathbf{x}, \mathbf{y})$ of $(R_{\mathbf{x}^*}, C_{\mathbf{x}^*})$ satisfies:
\[
  \mathcal{R}(\mathbf{x}, \mathbf{y}) > \frac12,\;\; \mathcal{C}(\mathbf{x}, \mathbf{y}) > \frac12
\]
The proof is by contradiction:

\begin{itemize}
  \item If the row player's payoff $\le 1/2$: since $\mathbf{x}^*$ is a maximin strategy, $\mathcal{R}(\mathbf{x}^*, \mathbf{y}) \ge \mathcal{R}(\mathbf{x}^*, \mathbf{y}^*) > 1/2$, so he can switch to $\mathbf{x}^*$ to improve his payoff, a contradiction;

  \item If the column player's payoff $\le 1/2$:

  Consider the zero-sum game $(-C_{\mathbf{x}^*}, C_{\mathbf{x}^*})$, and let $(\tilde{\mathbf{x}}, \tilde{\mathbf{y}})$ be a NE of it. Suppose $\mathcal{C}(\tilde{\mathbf{x}}, \tilde{\mathbf{y}}) \le 1/2$; then
  \[
    \mathcal{C}(\tilde{\mathbf{x}}, e_j) \le \mathcal{C}(\tilde{\mathbf{x}}, \tilde{\mathbf{y}}) \le 1/2 \;\; \text{for all } j \in [n]
  \]
  Moreover, $\operatorname{supp}(\tilde{\mathbf{x}}) \subseteq \operatorname{supp}(\mathbf{x}^*)$, so there exists $\tilde{\mathbf{x}}$ with $\operatorname{supp}(\tilde{\mathbf{x}}) \subseteq \operatorname{supp}(\mathbf{x}^*)$ and $\mathcal{C}(\tilde{\mathbf{x}}, e_j) \le \mathcal{C}(\tilde{\mathbf{x}}, \tilde{\mathbf{y}}) \le 1/2 \;\; \text{for all } j \in [n]$, which satisfies the condition of Case 1(b), a contradiction. Hence $\mathcal{C}(\tilde{\mathbf{x}}, \tilde{\mathbf{y}}) > 1/2$.

  Since $\tilde{\mathbf{y}}$ is the column player's maximin strategy in the zero-sum game, for every $\mathbf{x}$,
  \[
    \mathcal{C}(\mathbf{x}, \tilde{\mathbf{y}}) \ge \mathcal{C}(\tilde{\mathbf{x}}, \tilde{\mathbf{y}}) > \frac{1}{2}
  \]
  Then the column player can improve his payoff by switching to $\tilde{\mathbf{y}}$, contradicting that $(\mathbf{x}, \mathbf{y})$ is a NE.
\end{itemize}

Hence $(\mathbf{x}, \mathbf{y})$ is a $1/2$-WSNE of the original game. So we only need to compute such an $(\mathbf{x}, \mathbf{y})$. However, we cannot do this in polynomial time, because computing a NE of a general bimatrix game is PPAD-complete \cite{DGP09}.

Here we can use a corollary of the LMM lemma:

\begin{lemmaBox}{Uniform Sampling Lemma (Corollary of LMM [LMM03])}
  For any Nash equilibrium $(\mathbf{x}, \mathbf{y})$ of the restricted game $(R_{\mathbf{x}^*}, C_{\mathbf{x}^*})$, there exists a $\kappa(\delta)$-uniform strategy profile $(\mathbf{w}, \mathbf{z})$ such that:
  \[
    \forall i \in \operatorname{supp}(\mathbf{w}),\; \mathcal{R}(e_i, \mathbf{z}) \ge \mathcal{R}(\mathbf{x}, \mathbf{y}) - \delta > \frac12 - \delta
  \]
  \[
    \forall j \in \operatorname{supp}(\mathbf{z}),\; \mathcal{C}(\mathbf{w}, e_j) \ge \mathcal{C}(\mathbf{x}, \mathbf{y}) - \delta > \frac12 - \delta
  \]
  Since the maximum possible payoff is $1$ (the payoff matrices are normalized), the above conditions are equivalent to:
  \[
    \forall i \in \operatorname{supp}(\mathbf{w}),\ \mathcal{R}(e_i, \mathbf{z}) \ge \max_k \mathcal{R}(e_k, \mathbf{z}) - \left(1 - \left(\frac12 - \delta\right)\right) = \max_k \mathcal{R}(e_k, \mathbf{z}) - \left(\frac12 + \delta\right)
  \]
  and similarly for the column player. Therefore, $(\mathbf{w}, \mathbf{z})$ is a $(1/2+\delta)$-WSNE.
\end{lemmaBox}

Here $\kappa(\delta) = \left\lceil \dfrac{2 \ln (1/\delta)}{\delta^2} \right\rceil$ is a constant. A $k$-uniform strategy profile $(\mathbf{w}, \mathbf{z})$ satisfies that every component of $\mathbf{w}, \mathbf{z}$ is a non-negative integer multiple of $\dfrac{1}{k}$, so it naturally follows that $\max(|\operatorname{supp}(\mathbf{w})|, |\operatorname{supp}(\mathbf{z})|) \le \kappa(\delta)$, i.e., both supports have size at most $\kappa(\delta)$. Hence the total number of $\kappa(\delta)$-uniform strategies is $O(n^{\kappa(\delta)})$. Since the lemma guarantees existence, it suffices to enumerate all strategies and check whether each one is a $(1/2+\delta)$-WSNE. When $\delta$ is a constant, the running time of the algorithm is polynomial in $n$.

\begin{remarkBox}
  The difference from the LMM lemma here is that $\kappa(\delta)$ depends only on $\delta$, not on the problem size $n$. This is (probably) because in this case we only need to constrain the sampled support, whereas the LMM lemma needs to constrain all $n$ pure strategies.
\end{remarkBox}

\section{Algorithm 2: Query-Efficient \texorpdfstring{$(1/2+3\varepsilon+\delta)$}{1/2+3ε+δ}-WSNE}

\subsection{Motivation}
In practice, the cost of providing the entire payoff matrix may be unacceptable; only the payoff of a pure strategy profile $(e_i, e_j)$ can be queried. In this setting, the goal of the algorithm is to output an approximate WSNE with as few queries as possible.

If we keep using Algorithm 1, we face two problems:
\begin{enumerate}
  \item The preprocessing requires exact Nash equilibria of zero-sum games, but without access to the full matrices $R, C$ we can only compute approximate equilibria;
  \item Cases 1(b) and 1(c) require the full information of the restricted game $(R_{\mathbf{x}^*}, C_{\mathbf{x}^*})$, but the support of $\mathbf{x}^*$ may be large, making the number of queries unacceptable.
\end{enumerate}

\subsection{Modifications}

\subsubsection*{1. Approximate equilibria of zero-sum games}

For $n \times n$ zero-sum bimatrix games, there is a randomized algorithm that, with $O\left(\frac{n \log n}{\varepsilon^2}\right)$ queries, outputs an $\varepsilon$-NE with success probability $1 - n^{-1/n}$ [Goldberg \& Roth 2013].

Fearnley and Savani extended this result to the $\varepsilon$-WSNE case: with $O\left(\frac{n \log n}{\varepsilon^4}\right)$ queries, one can output an $\varepsilon$-WSNE with success probability $\left( 1 - n^{-\frac{1}{8}} \right)\left( 1 - \frac{2}{n} \right)^2$.

Now $(\mathbf{x}^*, \mathbf{y}^*)$ and $(\hat{\mathbf{x}}, \hat{\mathbf{y}})$ are no longer exact equilibria, and the relevant inequalities change accordingly.

By the definition of $\varepsilon$-WSNE:
\begin{align*}
  \mathcal{R}(\mathbf{x}^*, \mathbf{y}^*) &\ge \mathcal{R}(e_i, \mathbf{y}^*) - \varepsilon && \text{for all } i \in [n] \\
  \mathcal{C}(\hat{\mathbf{x}}, \hat{\mathbf{y}}) &\ge \mathcal{C}(\hat{\mathbf{x}}, e_j) - \varepsilon && \text{for all } j \in [n]
\end{align*}
Also, consider the corresponding form of the second inequality in the first game:
\begin{align*}
  -\mathcal{R}(\mathbf{x}^*, \mathbf{y}^*) &\ge -\mathcal{R}(\mathbf{x}^*, e_j) - \varepsilon && \text{for all } j \in [n] \\
  \iff \mathcal{R}(\mathbf{x}^*, \mathbf{y}^*) - \varepsilon &\le \mathcal{R}(\mathbf{x}^*, e_j) && \text{for all } j \in [n]
\end{align*}
Similarly, we have:
\[
  \mathcal{C}(\hat{\mathbf{x}}, \hat{\mathbf{y}}) - \varepsilon \le \mathcal{C}(e_i, \hat{\mathbf{y}}) \;\; \text{for all } i \in [n]
\]
Thus we obtain four inequalities for the approximate-NE setting. In the following, we call these four inequalities Observation 1 (in the original paper, this is called Observation 2).

\subsubsection*{2. Shrinking the support via the LMM lemma}

For the places in Algorithm 1 that need supports, taking $\mathbf{x}^*$ as an example, we sample a strategy $\mathbf{x}_s^*$ from $\mathbf{x}^*$ to replace it. Since the support of $\mathbf{x}_s^*$ has size $O(\log n / \epsilon^2)$, we can \textbf{query all} pairs $(e_i, e_j)$ with $i \in \operatorname{supp}(\mathbf{x}_s^*)$ and $j \in [n]$; the total number of such queries is $O(n \cdot \frac{\log n}{\epsilon^2})$.

Note that, to guarantee the correctness of the algorithm, the sampling cannot be arbitrary; it must satisfy the condition of Lemma 1 below.

The paper gives two lemmas:

\begin{lemmaBox}{Lemma 1 (Implied from the LMM lemma [LMM03])}
  Let $(\mathbf{x}, \mathbf{y})$ be a strategy profile of an $n \times n$ bimatrix game. Then there exists an $O(\log n / \varepsilon^2)$-uniform strategy profile $(x_s, y_s)$ sampled from $(\mathbf{x}, \mathbf{y})$ such that
  \begin{align*}
    |\mathcal{R}(\mathbf{x}_s, \mathbf{y}_s) - \mathcal{R}(\mathbf{x}, \mathbf{y})| &\le \varepsilon \\
    |\mathcal{C}(\mathbf{x}_s, \mathbf{y}_s) - \mathcal{C}(\mathbf{x}, \mathbf{y})| &\le \varepsilon \\
    |\mathcal{R}(e_i, \mathbf{y}_s) - \mathcal{R}(e_i, \mathbf{y})| &\le \varepsilon && \text{for all } i \in [n] \\
    |\mathcal{C}(\mathbf{x}_s, e_j) - \mathcal{C}(\mathbf{x}, e_j)| &\le \varepsilon && \text{for all } j \in [n]
  \end{align*}
\end{lemmaBox}

``Sampled from $(\mathbf{x}, \mathbf{y})$'' means that the sampling only considers the support of $(\mathbf{x}, \mathbf{y})$, ignoring all pure strategies with probability $0$.

Following the proof of the LMM lemma, for the last two conditions the failure probability is $O(n \exp(-2k\varepsilon^2))$; for the first two conditions, consider the decomposition:
\[
  |\mathcal{R}(\mathbf{x}_s, \mathbf{y}_s) - \mathcal{R}(\mathbf{x}, \mathbf{y})| \le |\mathcal{R}(\mathbf{x}_s, \mathbf{y}) - \mathcal{R}(\mathbf{x}, \mathbf{y})| + \max_{i} |\mathcal{R}(e_i, \mathbf{y}_s) - \mathcal{R}(e_i, \mathbf{y})|
\]
For the $\max$ term, the error probability is given by the union bound; for the first term, Hoeffding's inequality applies, giving failure probability $O(n \exp(-k\varepsilon^2))$.

With $k = c\ln n / \varepsilon^2$, the probability that the sampling fails the condition of Lemma 1 is a constant; increasing the constant $c$ in the number of samples makes this error probability decay exponentially.

\begin{lemmaBox}{Lemma 2}
  Let $(\mathbf{x}, \mathbf{y})$ be an $\varepsilon$-WSNE of an $n \times n$ bimatrix game, and let $(\mathbf{x}_s, \mathbf{y}_s)$ be an $O(\log n / \varepsilon^2)$-uniform strategy profile \textbf{satisfying the condition of Lemma 1}. Then $(\mathbf{x}_s, \mathbf{y}_s)$ is a $3\varepsilon$-WSNE of the game.

  \textbf{Proof:} Since Lemma 1 holds, for $i \in \operatorname{supp}(\mathbf{x}_s) \subseteq \operatorname{supp}(\mathbf{x})$ we have $|\mathcal{R}(e_i, \mathbf{y}_s) - \mathcal{R}(e_i, \mathbf{y})| \le \varepsilon$, hence
  \[
    \mathcal{R}(e_i, \mathbf{y}_s) \ge \mathcal{R}(e_i, \mathbf{y}) - \varepsilon \ge \max_{k \in [n]} \mathcal{R}(e_k, \mathbf{y}) - 2 \varepsilon
  \]
  because $(\mathbf{x}, \mathbf{y})$ is an $\varepsilon$-WSNE. What we need to prove is
  \[
    \mathcal{R}(e_i, \mathbf{y}_s) \ge \max_{k \in [n]} \mathcal{R}(e_k, \mathbf{y}_s) - 3 \varepsilon
  \]
  Let $k^*$ be the index attaining the maximum on the right-hand side; it suffices to prove $\mathcal{R}(e_i, \mathbf{y}_s) \ge \mathcal{R}(e_{k^*}, \mathbf{y}_s) - 3 \varepsilon$. Reusing Lemma 1:
  \begin{align*}
    &\mathcal{R}(e_{k^*}, \mathbf{y}_s) \le \mathcal{R}(e_{k^*}, \mathbf{y}) + \varepsilon \\
    \iff &\mathcal{R}(e_{k^*}, \mathbf{y}_s) - 3 \varepsilon \le \mathcal{R}(e_{k^*}, \mathbf{y}) - 2 \varepsilon \le \max_{k \in [n]} \mathcal{R}(e_k, \mathbf{y}) - 2 \varepsilon
  \end{align*}
  Together with the first inequality obtained above, the claim follows.
\end{lemmaBox}

It is worth noting that the statement of Lemma 2 in the paper does not mention the Lemma 1 restriction on the uniform strategy, and the correctness analysis of the subsequent algorithm also does not account for the probability of a failed Lemma 1 sampling. The only interpretation is that the authors implicitly assume that the constant in the sample size $k$ is large enough, so that the failure probability of the Lemma 1 sampling is negligible.

\subsection{Preprocessing}

\begin{enumerate}
  \item With $O\left(\frac{n \log n}{\varepsilon^4}\right)$ queries, obtain two approximate NEs $(\mathbf{x}^*, \mathbf{y}^*)$ and $(\hat{\mathbf{x}}, \hat{\mathbf{y}})$;
  \item Sample $\mathbf{x}_s^*, \mathbf{y}_s^*$ satisfying the condition of Lemma 1, query all $(e_i, e_j)$ with $i \in \operatorname{supp}(\mathbf{x}_s^*)$, and construct the restricted game $(R_{\mathbf{x}_s^*}, C_{\mathbf{x}_s^*})$.
\end{enumerate}

Then, following Algorithm 1, we case-analyze, again taking $\mathcal{R}(\mathbf{x}^*,\mathbf{y}^*) \ge \mathcal{C}(\hat{\mathbf{x}},\hat{\mathbf{y}})$ as an example:

\subsection{The \texorpdfstring{$\mathcal{R}(\mathbf{x}^*, \mathbf{y}^*) \ge \mathcal{C}(\hat{\mathbf{x}}, \hat{\mathbf{y}})$}{R(x*,y*) >= C(xh,yh)} branch}

\paragraph{Case 2(a): $\mathcal{R}(\mathbf{x}^*,\mathbf{y}^*) \le 1/2$}
Simply output $(\hat{\mathbf{x}}, \mathbf{y}^*)$.

\begin{itemize}
  \item Row player: $1/2 \ge \mathcal{R}(\mathbf{x}^*, \mathbf{y}^*) \ge \mathcal{R}(e_i, \mathbf{y}^*) - \varepsilon$, so every strategy is a $(1/2 + \varepsilon)$-best response;
  \item Column player: $1/2 \ge \mathcal{R}(\mathbf{x}^*, \mathbf{y}^*) \ge \mathcal{C}(\hat{\mathbf{x}}, \hat{\mathbf{y}}) \ge \mathcal{C}(\hat{\mathbf{x}}, e_j) - \varepsilon$ (by Observation 1), so every strategy is a $(1/2 + \varepsilon)$-best response.
\end{itemize}

Hence $(\hat{\mathbf{x}}, \mathbf{y}^*)$ is a $(1/2 + \varepsilon)$-WSNE, and in particular a $(1/2 + 3\varepsilon + \delta)$-WSNE.

\paragraph{Case 2(b): ``one high, one low''}
In this case, $\mathcal{R}(\mathbf{x}^*,\mathbf{y}^*) > 1/2$, and there exists $\mathbf{x}'$ supported on $\operatorname{supp}(\mathbf{x}_s^*)$ with $\mathcal{C}(\mathbf{x}', e_j) \le 1/2$; the output is $(\mathbf{x}', \mathbf{y}_s^*)$.

We use the two lemmas to show that $(\mathbf{x}', \mathbf{y}_s^*)$ is a $(1/2 + 3\varepsilon)$-WSNE:

\begin{itemize}
  \item Row player: by Lemma 2, $(\mathbf{x}_s^*, \mathbf{y}_s^*)$ is a $3\varepsilon$-WSNE, and $\mathbf{x}'$ is supported on $\operatorname{supp}(\mathbf{x}_s^*)$, which means that every pure strategy in the support of $\mathbf{x}'$ is a $3\varepsilon$-best response;
  \item Column player: by construction, the column player's maximum payoff is at most $1/2$.
\end{itemize}

Hence $(\mathbf{x}', \mathbf{y}_s^*)$ is a $(1/2 + 3\varepsilon)$-WSNE.

(Bonus: in the original paper, the justification for the row player mistakenly cites Lemma 1 instead of Lemma 2; a small typo, \sout{although Lemma 2 is also built upon Lemma 1}.)

\paragraph{Case 2(c): ``both high''}
In this case, neither of the previous two conditions holds.

Consider the restricted game $(R_{\mathbf{x}_s^*}, C_{\mathbf{x}_s^*})$. One can show that every Nash equilibrium $(\mathbf{x}, \mathbf{y})$ of $(R_{\mathbf{x}_s^*}, C_{\mathbf{x}_s^*})$ satisfies:
\[
  \mathcal{R}(\mathbf{x}, \mathbf{y}) > \frac12 - 3\varepsilon,\;\; \mathcal{C}(\mathbf{x}, \mathbf{y}) > \frac12 - 3\varepsilon
\]
The proof is again by contradiction:

\begin{itemize}
  \item If the row player's payoff $\le 1/2 - 3\varepsilon$:

  Following the previous proof, we try to show that the strategy can be switched to $\mathbf{x}_s^*$, i.e., $\mathcal{R}(\mathbf{x}_s^*, \mathbf{y}) > \frac12 - 3\varepsilon$.

  Consider the zero-sum game $(R, -R)$, where $(\mathbf{x}^*, \mathbf{y}^*)$ is an $\varepsilon$-WSNE of it. By the fourth inequality of Lemma 1:
  \[
    |\mathcal{C}(\mathbf{x}_s^*, e_j) - \mathcal{C}(\mathbf{x}^*, e_j)| \le \varepsilon \;\; \text{for all } j \in [n]
  \]
  Substituting the column player's payoff $-R$:
  \[
    |(-\mathcal{R})(\mathbf{x}_s^*, e_j) - (-\mathcal{R})(\mathbf{x}^*, e_j)| \le \varepsilon \;\; \text{for all } j \in [n]
  \]
  \[
    |\mathcal{R}(\mathbf{x}^*, e_j) - \mathcal{R}(\mathbf{x}_s^*, e_j)| \le \varepsilon \;\; \text{for all } j \in [n]
  \]
  Hence for every mixed strategy $\mathbf{y}$, $|\mathcal{R}(\mathbf{x}^*, \mathbf{y}) - \mathcal{R}(\mathbf{x}_s^*, \mathbf{y})| \le \varepsilon$. Also, by Observation 1:
  \[
    \mathcal{R}(\mathbf{x}^*, \mathbf{y}^*) - \varepsilon \le \mathcal{R}(\mathbf{x}^*, e_j) \;\; \text{for all } j \in [n]
  \]
  Therefore, for every mixed strategy $\mathbf{y}$,
  \[
    \mathcal{R}(\mathbf{x}_s^*, \mathbf{y}) \ge \mathcal{R}(\mathbf{x}^*, \mathbf{y}) - \varepsilon \ge \mathcal{R}(\mathbf{x}^*, \mathbf{y}^*) - 2\varepsilon > \frac{1}{2} - 2\varepsilon
  \]
  So the row player can deviate to $\mathbf{x}_s^*$ to satisfy the condition, a contradiction;

  \item If the column player's payoff $\le 1/2 - 3\varepsilon$:

  First, we state a claim:

  \begin{lemmaBox}{Assertion 1}
    For the restricted game $(-C_{\mathbf{x}_s}, C_{\mathbf{x}_s})$, let $(\tilde{\mathbf{x}}, \tilde{\mathbf{y}})$ be a NE of it. Then $\mathcal{C}(\tilde{\mathbf{x}}, \tilde{\mathbf{y}}) > 1 / 2 - 3\varepsilon$.
  \end{lemmaBox}

  This is quite intuitive: suppose $\mathcal{C}(\tilde{\mathbf{x}}, \tilde{\mathbf{y}}) \le 1/2 - 3 \varepsilon \le 1/2$; then $\tilde{\mathbf{x}}$ is a feasible solution to the condition of Case 2(b) (since $\operatorname{supp}(\tilde{\mathbf{x}}) \subseteq \operatorname{supp}(\mathbf{x}_s^*)$), which contradicts the premise of entering Case 2(c).

  Since $(\tilde{\mathbf{x}}, \tilde{\mathbf{y}})$ is a NE, $\tilde{\mathbf{y}}$ is the column player's maximin strategy, meaning that for every row strategy $\mathbf{x}$:
  \[
    \mathcal{C}(\mathbf{x}, \tilde{\mathbf{y}}) \ge \mathcal{C}(\tilde{\mathbf{x}}, \tilde{\mathbf{y}}) > \frac{1}{2} - 3\varepsilon
  \]
  Therefore, the column player can deviate to $\tilde{\mathbf{y}}$ to improve his payoff, contradicting that $(\mathbf{x}, \mathbf{y})$ is a NE.
\end{itemize}

Enumerating all $\kappa(\delta)$-uniform strategy profiles, we find a $(1/2+3\epsilon+\delta)$-WSNE. The time complexity analysis is the same as in Algorithm 1.

\section{Conclusion}

Overall, this is very impressive. Although the original paper has some jumps and ambiguities in the proofs, it is still very elegant: it simplifies the cases that need to be enumerated via a concise case analysis, and then uses the uniform sampling theorem to bound the enumeration time at the polynomial level. It would be even better if the proof of Algorithm 2 could be more detailed.
